The proof is in Gauge Theory Gravity (Fact \claimref{synthesis.imported-gtg-dictionary}): the metric of the statement is the derived metric of the gravitational field, its variation is the variation of $\hb$ with $\Omega$ the torsion-free field of $\hb$, and the Ricci and Riemann tensors are the Ricci map $\mathcal{R}(a) = \partial_{b} \cdot \mathcal{R}(b \wedge a)$ and the curvature operator $\mathcal{R}(B)$ in the store's sign. The statement's mixed components $R^{a}{}_{b}$ (signature $(-,+,+,+)$, Ricci tensor positive on spheres) are the components $\gamma^{a} \cdot \mathcal{R}(\gamma_{b})$ of the Ricci map on an orthonormal physical frame, and $f$ is a polynomial in these, contracted with the flat product.

\emph{(i).} Write $f = \sum_{k \geq 2} f_{k}$ with $f_{k}$ homogeneous of degree $k$ in the Ricci map, and let $\mathcal{P}_{f}$ be the linear map $\partial f/\partial\mathcal{R}(a)$, of degree at least one in the Ricci map. Vary $\hb$ in $\int\det(\hb)^{-1}f$, with $\Omega$ following as the torsion-free field: the variation of the weight $\det(\hb)^{-1}$ and of the frame in which $f$ is evaluated gives terms proportional to $f_{k}$ or to its derivative at fixed frame components, both of degree $k \geq 2$; the variation of the Ricci map gives $\mathcal{P}_{f}$ contracted with $\delta\mathcal{R}(a)$, which consists of terms algebraic in $\delta\hb$ times the field strength and of $D$-derivatives of $\delta\Omega$, itself a first $D$-derivative of $\delta\hb$; integration by parts leaves second $D$-derivatives of $\mathcal{P}_{f}$ and products of $\mathcal{P}_{f}$ with the field strength. Where the Ricci map vanishes, $\mathcal{P}_{f}$ vanishes identically, so all its derivatives vanish, and every term of the field equations of $f$ vanishes; the Einstein part vanishes because the field is vacuum. The Schwarzschild field in the Painlev\'e-Gullstrand gauge, $\hb = 1 - \beta P$, $\beta = \sqrt{2M/r}$, has the store-sign curvature operator
\[
\mathcal{R}(B) = \frac{M}{2r^{3}}\bigl(B + 3\sigma_{r}B\sigma_{r}\bigr), \qquad \sigma_{r} = \gamma_{1}\gamma_{0},
\]
with eigenvalue $2M/r^{3}$ on $\sigma_{r}$ and $I\sigma_{r}$ and $-M/r^{3}$ on the four other planes, so its Ricci map vanishes (the eigenvalues on the three planes through each frame vector sum to zero) and its Kretschmann scalar is $2\langle\mathcal{R}(a \wedge b)\,\mathcal{R}(\partial_{b} \wedge \partial_{a})\rangle = 48M^{2}/r^{6}$. The horizon locus is $Q(\hb^{-1}(e_{t})) = 1 - \beta^{2} = 0$, $r = 2M$, where $(K \cdot D)K = \Omega(e_{t}) \times K = K/(4M)$: surface gravity $1/4M$. The radial eigenvalue is $2M/r^{3} = \sqrt{K/12}$, so the surface of (C1), $2M/r^{3} = 1$ in Planck units, is the level set on which the field strength's radial eigenvalue is one; where it is below one, in particular where the field strength vanishes, it is absent. The check computes all of this exactly from the gauge field (rotation-gauge field from the torsion condition) at five rational points of two masses.

\emph{(ii).} By (W) and clause (b) of Fact \claimref{synthesis.imported-wald-noether-charge-entropy}, $S = -2\pi\oint E^{abcd}\epsilon_{ab}\epsilon_{cd}\,dA$, where $\epsilon_{ab}$ is the unit bivector $\sigma$ of the cross-section's normal plane ($\sigma^{2} = 1$, $\epsilon_{ab}\epsilon^{ab} = -2$). In GTG the contraction $E^{abcd}\epsilon_{ab}\epsilon_{cd}$ is the derivative of the Lagrangian density along the curvature operator $\mathcal{S}_{\sigma}(B) = -\langle B\sigma\rangle\,\sigma$ (the chain rule, with the sign fixed by the Einstein term and confirmed against the component contraction): $\mathcal{S}_{\sigma}$ has Ricci map $-P_{\sigma}$, minus the projector onto the normal plane, and scalar $-2$, so for $\mathcal{L} = \mathcal{R}/16\pi G$ the contraction is $-1/(8\pi G)$ and $S = A/4G$. For $f$ the derivative is $\mathcal{P}_{f}$ contracted with $-P_{\sigma}$, of degree at least one in the Ricci map; it vanishes on the Ricci-flat horizon, so $S = A/4G$ for every $f$ (the check computes it for $R_{ab}R^{ab}$ and $R^{2}$ on Schwarzschild).

\emph{(iii).} Theorem \claimref{singularity.rate-discharge-hypothesis-iii} derives the temperature $1/8\pi M$ from (P), (R) and the geometry of the Euclidean section, whose period $2\pi/\kappa$ depends only on the gravitational field, so it is the same under every $L_{f}$ by (i). Theorem \claimref{scrambling.mode-count} counts modes on the surface of (C1) with the density of (C2) over the slice volume of the interior: again the field only. Proposition \claimref{synthesis.quarter-from-temperature} integrates $dS = dM/T$ along the Schwarzschild family to $4\pi M^{2} = A/4$, which by (ii) is also the Wald entropy under every $L_{f}$. So any derivation that uses only these inputs reaches the same conclusions for every $f$. The field equations differ: for $f = \beta R_{ab}R^{ab}$, $\mathcal{P}_{f} = 2\beta\mathcal{R}$, and the second $D$-derivatives of it in (i) make the equations with matter of fourth order, not $\mathcal{G}(a) + 8\pi G\,T(a) = 0$; a derivation of Einstein's equations from data shared by all $L_{f}$ would derive them in every $L_{f}$, which is false for $f \neq 0$.

\emph{(iv).} The Reissner-Nordstr\"om field in the Painlev\'e-Gullstrand gauge, $\beta = \sqrt{2M/r - Q^{2}/r^{2}}$, has the store-sign Ricci map $\mathcal{R}(\gamma_{0}) = -\tfrac{Q^{2}}{r^{4}}\gamma_{0}$, $\mathcal{R}(\gamma_{1}) = -\tfrac{Q^{2}}{r^{4}}\gamma_{1}$ on the radial plane and $+\tfrac{Q^{2}}{r^{4}}$ on the two transverse frame vectors: $R^{t}{}_{t} = R^{r}{}_{r} = -Q^{2}/r^{4}$, nonzero at the horizon $r_{+}$ (the check computes it at $r = r_{+} = \tfrac52$ and $r = \tfrac92$ for $M = \tfrac32$, $Q^{2} = \tfrac54$). For $f = \beta R_{ab}R^{ab} = \beta\,\partial_{a} \cdot \mathcal{R}(\mathcal{R}(a))$ the derivative along $\mathcal{S}_{\sigma}$ is $-2\beta\,\partial_{a} \cdot \mathcal{R}(P_{\sigma}(a)) = -2\beta\,(R^{t}{}_{t} + R^{r}{}_{r})$, so the Wald entropy gains $4\pi\beta\,(R^{t}{}_{t} + R^{r}{}_{r})A = -8\pi\beta Q^{2}A/r_{+}^{4} \neq 0$ there. In Jacobson's derivation the heat flux through a local Rindler horizon is $\int k \cdot T(k)$ and, through Einstein's equations, equals the focusing term $-k \cdot \mathcal{R}(k)$; the Clausius relation is imposed for every point and every null $k$, which is where the Ricci map is nonzero and where (by the computation just made) the Ricci-dependent part of a horizon entropy is visible. For the remark: (R) is a premise about stationary Killing horizons and (C1) about Planck-curvature surfaces; by (i) a black hole's (C1) surface is absent from any region with $\sqrt{K/12} < 1$. The remark is a reading of the store, not part of the proved clauses.

\emph{In components} (a reading, signature $(-,+,+,+)$): the field equations of $\sqrt{-g}\,f$ are obtained by varying $g^{ab}$, with $P^{ab} = \partial f/\partial R_{ab}$ and $\delta R_{ab}$ second covariant derivatives of $\delta g$; the Christoffel computation of $g = \operatorname{diag}(-h, 1/h, r^{2}/(1 - x^{2}), r^{2}(1 - x^{2}))$ gives $R_{ab} = 0$, $K = 48M^{2}/r^{6}$ for $h = 1 - 2M/r$ and $R^{t}{}_{t} = R^{r}{}_{r} = -Q^{2}/r^{4}$ for $h = 1 - 2M/r + Q^{2}/r^{2}$; and $E^{abcd}\epsilon_{ab}\epsilon_{cd} = 2\beta R^{bd}\epsilon_{ab}\epsilon^{a}{}_{d} = -2\beta(R^{t}{}_{t} + R^{r}{}_{r})$ with $\epsilon_{tr} = -\epsilon_{rt} = 1$. The earlier check computes these, and the surface gravity $\tfrac12 h'(2M) = 1/4M$ and Hawking's integral along the family.
